\chapter*{General Introduction}
\addcontentsline{toc}{chapter}{General Introduction}
Performance evaluation becomes difficult when goals, progress evidence, comments, and HR decisions are scattered across documents, spreadsheets, private messages, and late review meetings. A manager may know that an employee worked hard, but the organization needs a traceable answer to more precise questions: what objectives were agreed, which KPIs were used, who approved them, how much progress was recorded, whether the employee had a chance to assess their own work, and how HR validated the final decision. The PerTrack project addresses that operational problem by turning a sensitive annual review process into a structured web application.

\begin{figure}[H]\centering
\resizebox{0.92\textwidth}{!}{%
\begin{tikzpicture}[every node/.style={font=\small}, stage/.style={draw, rounded corners, fill=PerTrackLight, minimum width=3.0cm, minimum height=0.95cm, align=center}, arr/.style={-Latex, thick}]
\node[stage] (setup) {Organization setup\\users, roles, teams};
\node[stage, right=1.0cm of setup] (phase1) {Phase 1\\objectives and KPIs};
\node[stage, right=1.0cm of phase1] (phase2) {Phase 2\\tasks and check-ins};
\node[stage, below=1.0cm of phase2] (phase3) {Phase 3\\evaluation};
\node[stage, left=1.0cm of phase3] (hr) {HR review\\validation or return};
\node[stage, left=1.0cm of hr] (follow) {Follow-up\\career, PIP, bonus};
\draw[arr] (setup) -- (phase1);
\draw[arr] (phase1) -- (phase2);
\draw[arr] (phase2) -- (phase3);
\draw[arr] (phase3) -- (hr);
\draw[arr] (hr) -- (follow);
\draw[arr, dashed] (hr.north) to[bend left=18] node[above, font=\scriptsize]{return for correction} (phase3.south west);
\end{tikzpicture}}
\caption{Performance-management lifecycle implemented by the application. Source: author's own work based on cycle phases and business modules.}
\label{fig:performance-lifecycle}
\end{figure}



The lifecycle in \cref{fig:performance-lifecycle} is the backbone of the report. It shows why the application cannot be reduced to isolated CRUD screens: each module prepares evidence for the next step. A cycle creates time boundaries, objectives define the performance contract, tasks and check-ins capture execution, evaluations transform evidence into judgment, and HR review gives the process institutional control.

The repository shows a system built around four roles: \texttt{ADMIN}, \texttt{HR}, \texttt{TEAM\_LEADER}, and \texttt{COLLABORATOR}. These roles are not decorative labels. They control route access in the backend, menu access in the frontend, and the visibility of objectives, evaluations, HR decisions, and audit information. This is important because performance data is personal and organizationally sensitive. A collaborator should be able to draft objectives and follow their own evaluation, while a team leader should review managed employees, and HR should validate the final process without bypassing every manager-owned step.

The application also gives special attention to objective weights. In this project, weights are not merely display values. They define how objective achievement contributes to the automatic evaluation score. A team objective is not divided by the number of members; instead, the full weight consumes capacity for every assigned member. This design is visible in both the scoring service and its tests. It prevents a shared objective from becoming artificially small when many employees participate, while still requiring the system to control the total allocation per employee and per team bucket.

PerTrack combines a React and Vite frontend with an Express and MongoDB backend. The frontend is organized around protected routes, role-based menus, dashboards, modal workflows, charts, kanban task views, active-cycle awareness, and token refresh logic. The backend exposes REST endpoints for users, teams, cycles, objectives, check-ins, tasks, meetings, feedback, evaluations, final evaluations, HR decisions, improvement plans, career recommendations, reports, notifications, calendar integration, audit logs, and AI assistance. MongoDB is used through Mongoose schemas with references and embedded arrays, which fits the project because objectives contain KPIs, comments, activity logs, change requests, and attachments that evolve during the review lifecycle \cite{mongodbDocs,mongooseSchemas}.

The project includes two forms of AI support. The first is generative assistance in the Node.js backend: goal suggestions, KPI suggestions, objective refinement, check-in drafts, mid-year review drafts, final self-assessment drafts, manager review drafts, and development-plan suggestions. The implementation is provider-aware and can use Gemini, xAI/Grok, or OpenAI-compatible chat completions depending on environment variables. The second is a Python Flask microservice trained on a synthetic dataset for performance rating and promotion-readiness prediction. That AI service loads saved XGBoost model artifacts, validates numeric inputs, computes a deterministic overall score, returns a predicted rating and promotion probability, and generates textual strengths, weaknesses, summaries, and suggestions. Because the dataset is synthetic, the report treats model metrics as project evidence, not as proof of real-world HR validity.

The engineering objective of the project is therefore wider than a simple CRUD application. The system must coordinate roles, workflow phases, objective approval, team hierarchy, self-assessment, manager scoring, HR review, auditability, and deployment. The implemented DevOps assets confirm this ambition: Dockerfiles for backend and frontend, a Compose file for local containers, Kubernetes base manifests, Kustomize overlays for dev, QA, staging, and production, an Argo CD application, and a GitHub Actions workflow with secret scanning, backend tests, frontend build/tests, Docker image build and push, Trivy scans, GitOps tag update, and a health-check smoke test.

This report follows the project from context to implementation. Chapter 1 presents the project environment and stakeholders. Chapter 2 studies the business problem, project scope, risks, and planning. Chapter 3 defines requirements and architecture. Chapters 4 to 7 follow implementation sprints covering organization/access, objectives/KPIs/tasks, evaluation/HR workflows, and AI/analytics. Chapter 8 covers testing, DevOps, deployment, and quality evidence. The report ends with limitations and realistic future work.
